% 【综合论述】所属：第6章（由 chapters/revised/06-formal.tex \input 引入）
\minitopic{后验概率不是行动按钮}同一患病概率对不同决定具有不同意义。是否进行低成本复检、是否使用高副作用药物、是否隔离患者，需要比较误报、漏报、等待、可逆性和资源占用。期望效用把每项行动在各状态下的后果按概率加权，但效用数值包含价值判断：谁承担损失、长期后果如何折现、公平约束是否可被金钱替代。把这些判断写进矩阵比隐藏在“风险过高”一句话中更诚实，也使争议能够定位。

\minitopic{阈值来自后果结构而非知识定义}若漏诊损失远大于误报，采取进一步检查的阈值可以很低；若行动不可逆且伤害严重，执行阈值应更高。两个阈值之间还可能存在“继续收集信息”区域。知识不必与任一固定概率阈值相同：彩票命题概率极高却缺少个案联系，直接知觉可能没有可报告数值却构成知识。概率为行动提供可比较输入，知识则评价信念如何连接事实。将二者区分后，我们既可说主体知道$p$却因巨大损失仍应复核，也可在尚无知识时依高概率采取可逆防护。

\minitopic{信息价值取决于能否改变行动}一项新检测即使非常准确，若结果不会改变任何可行行动，其决策价值可能很低；一种中等准确、快速而廉价的检测若能把个案从等待区移入行动区，价值反而较高。计算信息价值要比较检测前最优行动与看到每种结果后的最优行动，再扣除成本和延迟。该框架还提醒我们，收集更多数据并非总是理性：调查可能错过救治窗口、增加隐私风险或耗尽有限资源。停止调查需要理由，但“还有信息可得”本身不是继续的充分理由。

\minitopic{稳健决策不依赖单一点估计}若多个合理模型给出不同后验和效用，决策者可比较最坏后果、遗憾值或跨模型均表现可接受的方案。稳健不是永远选择最保守行动；过度保守也会产生机会成本和系统性伤害。关键是公开哪些不确定性足以改变选择，并为模型更新预设触发条件。例如当基础率超过某范围、传感器漂移达到某值或新亚群误差出现时，原决策规则自动进入复核。这样，模型不确定性被转化为制度可执行的监测。

\begin{argumentbox}
从数字到行动使用五步链：写出备选行动；列出相关状态与后果；用概率表示当前不确定性；计算或比较行动表现；最后检查公平、权利、角色义务和模型稳健性等不能被期望值完全吸收的约束。任何一步缺失，都可能把精确计算用于错误问题。
\end{argumentbox}

\minitopic{概率沟通必须保留条件方向}“检测对患者有90\%概率阳性”与“阳性者有90\%概率患病”方向不同；“未来十年发生概率30\%”也不表示事件将在第三年发生。自然频率、绝对风险和对照基线通常比孤立百分比更易理解。沟通还要报告不确定范围、适用人群和版本时间。只给一个醒目数字会制造虚假确定性，使听者无法判断证据更新还是模型设定导致结论变化。形式素养不仅是会算，更是能把计算条件完整传递给他人。

\minitopic{模型责任需要分层}数据团队负责采样与测量，建模者负责假设和验证，领域专家负责解释适用性，决策者负责把概率与后果结合，机构负责监测和申诉。一次失败可能跨越多层：模型在原人群校准良好，却被部署到新环境；医生看见警告却没有时间复核；制度又没有记录覆写理由。把责任全部归给“算法”或最后点击按钮的人都会遗漏因果与规范结构。可问责系统应让每层留下可检查产物，并明确谁能暂停使用。

\minitopic{效用数字不会消除价值分歧}把后果换算成同一尺度便于比较，却可能把不可替代的生命、权利与尊严压成一个方便数字。即使接受量化，也要区分个人内不同结果的权衡与不同人之间的聚合：一百人各获得很小便利，是否足以抵消一人承受不可逆伤害，并非概率公理能够回答。决策报告应分别展示概率判断、后果估值和聚合原则，使反对者能够指出争议究竟位于事实预测还是价值计算。隐藏价值不是中立，只是让默认权重不受审查。

\minitopic{预防原则也需要触发与退出条件}面对低概率而灾难性后果，决策者可能在证据尚弱时采取防护。预防并不等于永远选择最坏情形；若任何逻辑可能都触发最大干预，行动会互相冲突且资源迅速耗尽。合理方案要说明伤害严重度、可能性证据、措施可逆性与替代成本，还要规定何时降级或终止。这样，预防原则成为处理尾部风险的一种结构化判断，而不是在需要支持既定立场时才使用的口号。

\minitopic{权利与角色义务可以限制期望最大化}医生对患者、法官对被告、公共机构对公民承担的义务，不总能被总体效用吸收。即使错误拘押少数人能提高平均安全感，正当程序仍可能禁止这种方案；即使公开个人信息能改善模型，也需满足同意与用途限制。决策模型可以把权利写成不可越过的约束、词典式优先级或极高代价，但采用哪种方式本身是规范选择。形式清晰的好处不是证明约束正确，而是显示哪些行动因何被排除。

\minitopic{群体决定需要说明谁的概率进入系统}委员会成员可能掌握不同证据、使用不同模型，也可能在讨论后彼此锚定。简单平均置信度忽略专长与证据相关性，按权威加权又可能压低异议。更好的程序先汇集各自理由和来源，再判断独立性与专长范围，必要时保留少数报告。群体最后采取一项行动，不表示成员必须拥有同一信念；重要的是决定规则能保存不确定、记录反对并在新证据到来时重新开启。

\minitopic{行动会反过来改变数据环境}预测高风险后采取干预，风险没有实现，不能直接说明模型错误；反之，资源按模型分配会使被忽视群体数据更少，未来模型看起来更不需要关注他们。信用、治安与医疗系统尤其容易形成反馈：模型影响行动，行动影响可观察结果，可观察结果再训练模型。验证必须区分自然结果与干预后结果，并寻找未部署对照、机制分析或审慎实验。否则系统会把自己的影响当成世界原本如此。

\minitopic{可逆与不可逆决定应采用不同证据策略}在可低成本撤回的行动中，可以较早试行并密切监测；一旦造成不可恢复伤害、锁定基础设施或改变权利地位，进入门槛应更高，试验范围应更小。这个区分把认识不足与行动设计连接起来：证据弱并不总要求完全等待，也可要求分阶段、可回滚和预设停止线。相同后验概率因行动可逆性不同而支持不同策略，不是概率标准任意变化，而是后果结构改变。

\minitopic{复盘必须检验当时过程而非只看结果}好决定可能遇到坏结果，坏决定也可能侥幸成功。复盘若只奖惩结果，会鼓励后见偏差与风险掩盖。应保存当时可得数据、模型版本、置信区间、替代方案和价值权重，再问过程是否按证据强度和既定规则运行。随后才用结果更新模型：哪些似然估计错了，哪个故障未列入，哪项代价被低估。过程评价与结果学习结合，才能既不因运气错责，也不因主观善意停止改进。

\begin{argumentbox}
一项负责任决定应同时留下概率档案与规范档案。前者记录数据、模型、依赖和敏感性，后者记录受影响者、后果权重、权利约束、决定权限、回滚条件与申诉入口。只有两份档案能够相互对应，数字才真正进入了可问责行动。
\end{argumentbox}

\minitopic{阶段性结论}形式认识论的价值不是把哲学问题消解成算术，而是迫使我们写清假设、证据方向、依赖结构和价值权衡。概率一致性限制信念组合，贝叶斯更新说明新证据如何改变置信，决策理论比较不确定中的行动；但问题选择、模型适用、知识资格和制度责任仍需更广认识论。最成熟的使用方式，是让数字增加可检验性，同时拒绝让数字遮蔽它无法回答的问题。
